% =====================================================================
% SWAT+ ENGINE ACCELERATION — Supplementary Material (Computers & Geosciences).
% Compiled separately from the main manuscript; shares the preamble and bib.
%   pdflatex supplementary && bibtex supplementary && pdflatex supplementary && pdflatex supplementary
% =====================================================================
\documentclass[12pt]{article}
% Shared preamble — SWAT+ engine-acceleration manuscript.
\usepackage[margin=1in]{geometry}
\usepackage{setspace}
\usepackage{graphicx}
\usepackage{booktabs}
\usepackage{array}
\usepackage{tabularx}
\usepackage{ragged2e}
\usepackage{longtable}
\usepackage{natbib}
\usepackage{amsmath}
\usepackage{listings}
\usepackage{xcolor}
\usepackage[hidelinks]{hyperref}
\usepackage{url}

\DeclareRobustCommand{\path}[1]{\nolinkurl{#1}}

\newcolumntype{Y}{>{\RaggedRight\arraybackslash}X}

\graphicspath{{../figures/}}

% Lightweight code listing style for the Fortran snippets (threadprivate, wavefront).
\lstdefinestyle{swxcode}{
  basicstyle=\ttfamily\small,
  keywordstyle=\color{blue!60!black},
  commentstyle=\color{green!40!black}\itshape,
  breaklines=true,
  frame=single,
  framesep=4pt,
  showstringspaces=false,
  language=Fortran,
}
\lstset{style=swxcode}

% Inline TODO marker for the drafting phase (strip before submission).
\newcommand{\TODO}[1]{\textcolor{red}{[\textbf{TODO:} #1]}}

% Headline parallel/total numbers (final engine, dedicated c8a AWS instance; consistent re-run).
\newcommand{\parHead}{5.33}    % peak self-relative thread speedup (c8a, N=24, still climbing)
\newcommand{\parHeadN}{24}     % thread count at that peak
\newcommand{\totHead}{7.14}    % combined stock->parallel total on c8a at 24 threads
\newcommand{\totMinutes}{half a minute} % per-year wall after parallelization (27 s on c8a)

% Figure placeholder: keeps the manuscript compiling before the image files exist.
% Swap \figplaceholder{...} for \includegraphics{...} once the figure is generated.
\newcommand{\figplaceholder}[2]{%
  \fbox{\begin{minipage}[c][3.2cm][c]{0.88\textwidth}\centering
    \textcolor{red}{\textbf{[FIGURE TO GENERATE: \texttt{#1}]}}\\[4pt]
    \textcolor{red}{\footnotesize #2}
  \end{minipage}}%
}


% Supplementary numbering: S1, S2, ...; figures Figure_S1, tables Table_S1.
\renewcommand{\thesection}{S\arabic{section}}
\renewcommand{\thefigure}{S\arabic{figure}}
\renewcommand{\thetable}{S\arabic{table}}
\renewcommand{\theequation}{S\arabic{equation}}

\title{Supplementary Material\\[2pt]
\large Accelerating a Regional Hyper-Resolution {SWAT+} Model Without Changing the Science}
\author{Vahid Rafiei}
\date{}

\begin{document}
\maketitle

This document provides reproducibility detail, the Fortran reentrancy hazard and its code,
the full correctness forensics, and the extended scaling tables that support the main text.
All measurements use the same benchmark model and protocol described in the main paper (Methods, Section~3). Section, figure, and table numbers prefixed with ``S'' refer to this document.

% =====================================================================
\section{Reproducibility: per-optimization commits and the build recipe}\label{si:repro}
Every optimization in the cumulative ladder is a single commit in the engine fork
(\path{https://github.com/rafiei-vahid/swatplus}), so the marginal contribution of each rung is
auditable commit by commit (Table~\ref{tab:commits}). The benchmark harness, the measured
result files, the figure-generation script, and the benchmark model are released at
\path{https://github.com/SWATGenX/swatplus-engine-acceleration}.

\begin{table}[h]
\centering
\caption{The cumulative-ladder rungs, each a single commit. ``Class'' is the bottleneck the
change addresses. Rungs R1/R2---the NetCDF output backend and the gauge print filter---are the I/O
extensions documented and benchmarked in the companion paper \citep{Rafiei2026SWATGenX}; they are
available within R0 here and not re-derived, so the compute ladder begins at R3.}
\label{tab:commits}
\begin{tabular}{lll}
\toprule
Rung & Commit & Optimization (class) \\
\midrule
R0  & \texttt{768f1d1} & Stock engine baseline (NetCDF/print-filter I/O backend available) \\
R3  & \texttt{87abf8e} & \texttt{hru\_read} $\mathcal{O}(1)$ name index (startup) \\
R4  & \texttt{247e95b} & \texttt{varinit} per-row reset (per-step) \\
R5b & \texttt{91922e3} & Drop redundant inflow-hydrograph struct copy (channel) \\
R5c & \texttt{6d41300} & \texttt{ch\_temp} reset scope (channel) \\
R6  & \texttt{34f5c02} & Narrow-level fusion; final engine (serial + OpenMP builds) \\
\bottomrule
\end{tabular}
\end{table}

\paragraph{Toolchain.} The engine builds with the Intel oneAPI Fortran compiler (\texttt{ifx}) at
\texttt{-O3 -ipo -recursive -init=zero -init=arrays -free -fpe0}; the OpenMP build adds
\texttt{-fiopenmp} (\texttt{-DSWATPLUS\_OPENMP=ON}). \texttt{gfortran} fails to produce a working
binary at this model size. Two build practicalities are easy to miss: a fresh CMake configure must
pin \texttt{-DCMAKE\_Fortran\_COMPILER=ifx} \emph{and}
\texttt{PKG\_CONFIG\_PATH=<deps>/netcdf-ifx/lib/pkgconfig} (otherwise CMake selects
\texttt{gfortran}/system NetCDF and the \texttt{.mod} files clash), and the Intel \texttt{setvars.sh}
must be sourced without \texttt{set -u}. The \texttt{ifx} binaries link the Intel runtime and
NetCDF (\texttt{libnetcdf.so.19}, \texttt{libnetcdff.so.7}, \texttt{libhdf5}); on a stock Ubuntu
host these are \texttt{apt}-installable.

\paragraph{One binary, two modes.} The final OpenMP binary runs full-wavefront by default and the
byte-identical land-parallel/routing-serial mode when \texttt{SWATPLUS\_ROUTING\_SERIAL=1} is set
in the environment (Section~\ref{si:forensics}).

% =====================================================================
\section{The Fortran reentrancy hazard}\label{si:reentrancy}
Making the per-object body reentrant required auditing every module-global ``current-object''
scratch variable and giving each thread its own copy via OpenMP \texttt{threadprivate} (main text, Section~3.4). A Fortran-specific hazard complicates the na\"ive fix. A local declared with an
initializer,
\begin{verbatim}
      integer :: j = 0   ! implicitly SAVE in Fortran -> shared across calls
                         ! -> shared across threads -> a race
\end{verbatim}
is \emph{implicitly} \texttt{SAVE}d, hence shared across calls and therefore across threads. The
reflexive fix---stripping the initializer so the local becomes automatic (per-thread under
\texttt{-recursive})---is correct \emph{only} for write-before-read scratch. For the small set of
accumulators that legitimately persist across calls, stripping the \texttt{SAVE} silently broke
cross-call state and changed results \emph{even at one thread}. The correct treatment is to keep
those few variables \texttt{SAVE}d \emph{and} mark them \texttt{threadprivate}, so each thread
retains its own persistent copy. This was caught by the one-thread byte-identity check
(Section~\ref{si:forensics}): it surfaced a small ($\sim$0.3\%) \emph{deterministic} residual in
\texttt{ch\_watqual4} and \texttt{et\_pot} where stripping \texttt{SAVE} had dropped persistent
scratch; restoring it as \texttt{threadprivate}/\texttt{SAVE} state made one-thread output
bit-identical again.

% =====================================================================
\section{Correctness forensics}\label{si:forensics}
\paragraph{Thread-count invariance as the acceptance test.} For $N>1$ we require repeated runs at
different thread counts (1-vs-4, 4-vs-4, 4-vs-8) to produce identical daily output to a strict
tripwire ($>10^{-6}$ counts as a difference). Because a correct shared-memory program is
deterministic regardless of scheduling, any run-to-run difference at a fixed thread count is, by
definition, a data race, and the specific output fields that differ point at the unprotected state.
We separate genuine races from benign arithmetic with two builds: a non-\texttt{ipo} build,
bit-identical to the serial source (proving the algorithm is unchanged), and the \texttt{-ipo}
build whose residual ($\sim$$5\times10^{-4}$ relative) is floating-point reassociation.

\paragraph{The particulate in-stream race and its compiler root cause.} A streamflow-only one-year
test is insufficient: it passes while a latent race in the in-stream \emph{particulate}
constituents goes undetected. Only when in-stream water quality is exercised over a multi-year run
does thread-count invariance expose a divergence of order $15$--$20\%$ in three
species---organic-N, ammonia, and sediment-bound phosphorus---with run-to-run variation at higher
thread counts confirming a race rather than reassociation. The HRU land phase is byte-identical
across thread counts, localizing the race to the channel routing phase; the selectivity
(particulate-only, dissolved NO$_3$/soluble-P spared) localizes it further to the sediment-
associated transport. We traced it with ThreadSanitizer (\texttt{ifx -fsanitize=thread}; the
Archer-aware OpenMP runtime suppresses spurious synchronization reports). The race is \emph{not} a
missing \texttt{threadprivate}: every per-object index and scratch buffer is privatized. It is that
\texttt{ifx} allocates the \emph{temporaries} of the 152-byte \texttt{hyd\_output} derived type
(operator results, function returns, and argument copy-ins, used for every in-stream mass-balance
expression) in \textbf{shared static module storage} rather than on the per-thread stack;
concurrent channels then clobber the same static temporary mid-expression. The behavior is
invariant to every compiler control we tried (\texttt{-recursive}, \texttt{-auto},
\texttt{-heap-arrays}, \texttt{elemental} operators, and disabling zero-initialization). A targeted
source rewrite replacing the operator expressions with field-by-field helper subroutines is
byte-identical at one thread but does \emph{not} remove the race: passing a derived-type array
element as a subroutine argument re-creates a static copy-in temporary. A complete fix requires
eliminating every derived-type temporary in every concurrently executed routine---inline, with no
operators and no derived-type arguments---which is large and brittle, or a future compiler that
stack-allocates these temporaries.

\paragraph{Resolution.} Because the race lives entirely in the routing phase, running the land
phase in parallel and the routing phase serially is race-free and byte-identical (run-to-run
identical to $\sim$$10^{-12}$ on basin constituent totals; thread-count-invariant to the
floating-point-reassociation floor). This is the byte-identical mode of Table~\ref{tab:fallbackfull};
the full wavefront---faster, bit-exact for flow/sediment/dissolved species---is the alternative
mode for non-water-quality or tolerance-accepting use.

% =====================================================================
\section{Extended scaling tables}\label{si:scaling}
Table~\ref{tab:crosshwfull} gives the full cross-hardware scaling grid behind the main-text
figure, including the stock baseline and serial gain per machine. Table~\ref{tab:fallbackfull}
gives the full-wavefront versus byte-identical comparison on the reference node.

\begin{table}[h]
\centering
\caption{Cross-hardware scaling of the final engine (each machine's own $1\rightarrow N$ ratio),
identical engine, one-simulated-year benchmark workload, daily \texttt{channel\_sd} output,
thread-pinned, best of two, on dedicated AWS instances.}
\label{tab:crosshwfull}
\begin{tabular}{lccc}
\toprule
 & c8a (Turin) & c8i (Granite Rapids) & c5a (Rome) \\
\midrule
physical cores (vCPU) & 32 (32) & 16 (32) & 16 (32) \\
SMT                   & off & on ($2\times$) & on ($2\times$) \\
stock $R_0$ (s) & 191.9 & 272.2 & 350.6 \\
serial gain     & 1.67$\times$ & 2.47$\times$ & 2.16$\times$ \\
$t_1$ (s)       & 143.1 & 130.1 & 202.5 \\
\midrule
2  & 1.70$\times$ & 1.59$\times$ & 1.62$\times$ \\
4  & 2.79$\times$ & 2.41$\times$ & 2.47$\times$ \\
6  & 3.48$\times$ & 2.93$\times$ & 2.89$\times$ \\
8  & 3.95$\times$ & 3.29$\times$ & 3.21$\times$ \\
10 & 4.30$\times$ & 3.55$\times$ & 3.42$\times$ \\
12 & 4.54$\times$ & 3.75$\times$ & 3.60$\times$ \\
16 & 4.90$\times$ & \textbf{4.01$\times$} & \textbf{3.81$\times$} \\
20 & 5.18$\times$ & 3.73$\times$ & 3.51$\times$ \\
24 & \textbf{5.33$\times$} & 3.74$\times$ & 3.58$\times$ \\
\bottomrule
\end{tabular}
\end{table}

\begin{table}[h]
\centering
\caption{Full wavefront vs.\ byte-identical (routing-serial) mode on a dedicated c8a instance,
same final engine. Self-relative speedup and the forfeit (the parallel speedup the byte-identical
mode gives up by serializing routing).}
\label{tab:fallbackfull}
\begin{tabular}{rccc}
\toprule
Threads & Full wavefront & Byte-identical & Forfeit \\
\midrule
1  & 142.9\,s \quad 1.00$\times$ & 143.5\,s \quad 1.00$\times$ & --- \\
2  & 84.2\,s \quad 1.70$\times$  & 97.9\,s \quad 1.47$\times$  & 14\% \\
4  & 51.2\,s \quad 2.79$\times$  & 70.3\,s \quad 2.04$\times$  & 27\% \\
8  & 35.9\,s \quad 3.98$\times$  & 57.4\,s \quad 2.50$\times$  & 37\% \\
16 & 29.3\,s \quad 4.88$\times$  & 51.0\,s \quad 2.81$\times$  & 43\% \\
\bottomrule
\end{tabular}
\end{table}

\paragraph{A note on the runs.} The full-versus-byte-identical comparison (Table~\ref{tab:fallbackfull})
is an \emph{independent} run on the reference c8a node, with its own stock baseline
$R_0=193.5$\,s---within $0.8\%$ of the cross-hardware run's $191.9$\,s (Table~\ref{tab:crosshwfull}).
The three slightly different one-thread parallel walls ($143.1$\,s in Table~\ref{tab:crosshwfull};
$142.9$ and $143.5$\,s here) are these independent best-of-two measurements, not a discrepancy.

\paragraph{Composing the layers.} The end-to-end speedup is the product of three independent layers:
output scope (an I/O effect, quantified for these models in the companion benchmark
\citep{Rafiei2026SWATGenX} and excluded from the main figures), the machine-portable serial gain
($1.67$--$2.47\times$), and the thread scaling above. Against the comparison run's
$R_0=193.5$\,s the directly-measured stock$\rightarrow$parallel total is $3.78\times$ (full) /
$2.75\times$ (byte-identical) at four threads and $5.39\times$ / $3.37\times$ at eight; on the
cross-hardware run ($R_0=191.9$\,s) the full wavefront reaches $7.14\times$ at 24 threads.

\bibliographystyle{plainnat}
\bibliography{../bib/references-engine}
\end{document}
